\usepackage{morewrites} % <-- Adicionado o mais alto possível!

% !TeX TXS-program:compile = txs:///pdflatex | txs:///makeindex %.nlo -s nomencl.ist -o %.nls | txs:///pdflatex

\usepackage{morewrites} % <-- Importante estar no topo da lista de pacotes!

%-------------------------------------------------------
% CONFIGURAÇÃO DA PÁGINA
%-------------------------------------------------------
\usepackage[top=3cm, bottom=2cm, right=2cm, left=3cm]{geometry} % Margens ABNT

%-------------------------------------------------------
% IDIOMA E CODIFICAÇÃO
%-------------------------------------------------------
\usepackage[utf8]{inputenc}      % Permite digitação em UTF-8 (acentuação)
\usepackage[brazilian]{babel}    % Traduz automaticamente nomes para o português
\usepackage{silence}
\WarningFilter{tracklang}{No `datatool' support for dialect}

%-------------------------------------------------------
% PDF
%-------------------------------------------------------
\usepackage[final]{pdfpages}     % Insere páginas de arquivos PDF

%-------------------------------------------------------
% CITAÇÕES, LEGENDAS E NUMERAÇÃO
%-------------------------------------------------------
\usepackage{csquotes}            % Aspas inteligentes
\usepackage[labelsep=endash]{caption} % Formatação das legendas de figuras e tabelas

% Desvincula a numeração de figuras e tabelas dos capítulos (1, 2, 3...)
\counterwithout{figure}{chapter}
\counterwithout{table}{chapter}

%-------------------------------------------------------
% PARÁGRAFOS
%-------------------------------------------------------
\usepackage{indentfirst}         % Indenta o primeiro parágrafo de cada seção

%-------------------------------------------------------
% AMBIENTES PERSONALIZADOS (Gráficos e Quadros)
%-------------------------------------------------------
\usepackage{newfloat}            % Permite criar novos ambientes flutuantes

\DeclareFloatingEnvironment[
    name=Gráfico, 
    listname={Lista de Gráficos},
    within=none % Garante numeração contínua (1, 2, 3...)
]{grafico}

\DeclareFloatingEnvironment[
    fileext=loq,
    listname={Lista de Quadros},
    name=Quadro,
    within=none % Garante numeração contínua (1, 2, 3...)
]{quadro}

%-------------------------------------------------------
% TEXTO DE EXEMPLO E APÊNDICES
%-------------------------------------------------------
\usepackage{blindtext}           % Gera texto fictício para testes
\usepackage{lipsum}              % Gera texto Lorem Ipsum
\usepackage[titletoc]{appendix}  % Melhor controle dos apêndices

% Altera o título da bibliografia
\addto{\captionsbrazil}{%
\renewcommand{\bibname}{\normalsize\bfseries\hfill REFERÊNCIAS\hfill}
}

%-------------------------------------------------------
% CABEÇALHO E NUMERAÇÃO DE PÁGINAS
%-------------------------------------------------------
\usepackage{fancyhdr}
\setlength{\headheight}{14.5pt} 

% Configura o estilo padrão para todas as páginas
\pagestyle{fancy}
\fancyhf{} % Limpa configurações antigas de cabeçalho e rodapé
\fancyhead[R]{\thepage} % Número da página no topo (Right/Direita)
\renewcommand{\headrulewidth}{0pt} % Remove a linha horizontal abaixo do número

% Redefine o estilo "plain" (usado automaticamente no início dos capítulos)
\fancypagestyle{plain}{
    \fancyhf{}
    \fancyhead[R]{\thepage}
    \renewcommand{\headrulewidth}{0pt}
}

\usepackage{epigraph}            % Cria epígrafes
\renewcommand{\textflush}{flushepinormal} % Alinhamento do texto da epígrafe

%-------------------------------------------------------
% FIGURAS E TABELAS
%-------------------------------------------------------
\usepackage{graphicx}            % Inclusão de imagens
\usepackage{float}               % Posicionamento preciso de figuras e tabelas (H)
\usepackage{capt-of}             % Permite usar legendas fora de ambientes flutuantes
\usepackage{booktabs}            % Tabelas com melhor qualidade tipográfica
\usepackage{longtable}           % Necessário para o glossário customizado

%-------------------------------------------------------
% DESENHOS E DIAGRAMAS
%-------------------------------------------------------
\usepackage{tikz}                % Desenhos vetoriais
\usetikzlibrary{arrows.meta,bending} % Bibliotecas extras do TikZ
\usepackage[all]{xy}             % Diagramas comutativos (matemática)

%-------------------------------------------------------
% MATEMÁTICA E FÍSICA
%-------------------------------------------------------
\usepackage{amsmath}             % Fórmulas matemáticas
\usepackage{amssymb}             % Símbolos matemáticos
\usepackage{amsfonts}            % Fontes matemáticas
\usepackage{amsthm}              % Ambientes de teoremas
\usepackage{MnSymbol}            % Símbolos matemáticos adicionais
\usepackage{thmtools}            % Gerenciamento avançado de teoremas
\usepackage{physics}             % Comandos úteis para Física e Matemática
\usepackage{cancel}              % Cancela termos em equações

\DeclareMathOperator{\mdc}{mdc}  % Operador matemático "mdc"
\DeclareMathOperator{\sen}{sen}  % Operador matemático "sen"

%-------------------------------------------------------
% TEXTO E CAIXAS
%-------------------------------------------------------
\usepackage{hyphenat}            % Controle da hifenização
\usepackage{setspace}            % Controle do espaçamento entre linhas
\linespread{1.25}                % Espaçamento de 1,25
\usepackage{boxedminipage}       % Cria caixas envolvendo blocos de texto
\usepackage{enumitem}            % Personalização de listas
\usepackage{changepage}          % Permite alterar margens temporariamente

\newenvironment{citacao}
{\begin{quote}\small\singlespacing}
{\{end{quote}}

%-------------------------------------------------------
% NOMENCLATURA E SIGLAS
%-------------------------------------------------------
\usepackage{nomencl}             % Lista de símbolos (nomenclatura)
\makenomenclature                % Inicializa o gerador de símbolos
\setlength{\nomlabelwidth}{3cm}
\setlength{\nomitemsep}{-\parsep}
\renewcommand{\nomname}{\normalsize \makebox[\linewidth]{\MakeUppercase{Lista de Símbolos}}}

\usepackage[acronym,nogroupskip]{glossaries}
\makeglossaries                  % Inicializa o gerador de siglas
\renewcommand{\glossarysection}[2][]{}

\newglossarystyle{siglasABNT}{
  \setglossarystyle{long}
  \renewenvironment{theglossary}
  {\begin{longtable}{@{}p{3cm}p{12cm}@{}}}
  {\end{longtable}}
  \renewcommand*{\glossentry}[2]{%
    \glstarget{##1}{\textbf{\glossentryname{##1}}} & \glossentrydesc{##1} \\
  }
}

\newcommand{\sigla}[2]{%
    \newacronym{#1}{#1}{#2}%
    \gls{#1}%
}

% --- DEFINIÇÃO DAS SIGLAS DO TRABALHO ---
\newacronym{etl}{ETL}{\textit{Extract, Transform and Load}}
\newacronym{elt}{ELT}{\textit{Extract, Load, Transform}}
\newacronym{ssis}{SSIS}{SQL Server Integration Services}
\newacronym{dag}{DAG}{\textit{Directed Acyclic Graph}}
\newacronym{sefazpb}{SEFAZ-PB}{Secretaria de Estado da Fazenda da Paraíba}
\newacronym{atf}{ATF}{Sistema de Administração Tributária e Financeira}
\newacronym{sor}{SOR}{\textit{System of Record}}
\newacronym{scamf}{SCAMF}{Sistema de Controle e Análise das Malhas Fiscais}
\newacronym{nfe}{NFe}{Nota Fiscal Eletrônica}
\newacronym{nfce}{NFCe}{Nota Fiscal de Consumidor Eletrônica}
\newacronym{cte}{CTe}{Conhecimento de Transporte Eletrônico}
\newacronym{efd}{EFD}{Escrituração Fiscal Digital}
\newacronym{bi}{BI}{\textit{Business Intelligence}}
\newacronym{nutes}{Nutes}{Núcleo de Tecnologia Estratégica em Saúde}
\newacronym{ufpb}{UFPB}{Universidade Federal da Paraíba}

%-------------------------------------------------------
% DATAS E TÍTULOS DE SEÇÃO
%-------------------------------------------------------
\usepackage{datetime}
\usepackage{titlesec}

\titleformat*{\section}{\normalsize \bfseries}
\titleformat*{\subsection}{\normalsize\itshape\bfseries}
\titleformat*{\subsubsection}{\normalsize \bfseries}
\titleformat*{\paragraph}{\large\bfseries}
\titleformat*{\subparagraph}{\large\bfseries}
\titleformat{\chapter}{\bfseries}{\filright \thechapter}{20pt}{\filright}
\titlespacing*{\chapter}{0pt}{0pt}{*1.5}
\titlespacing{\section}{0pt}{*1.5}{*1.5}

%-------------------------------------------------------
% TOCLOFT: FORMATAÇÃO DO SUMÁRIO (ABNT NBR 6027: 2012)
%-------------------------------------------------------
\usepackage{tocloft}

\renewcommand\cftchapnumwidth{3.8em}
\renewcommand\cftsecnumwidth{3.8em}
\renewcommand\cftsubsecnumwidth{3.8em}
\renewcommand{\cftsecfont}{\bfseries }

\renewcommand{\cftfigpresnum}{Figura~}
\renewcommand{\cftfigaftersnum}{ --}
\settowidth{\cftfignumwidth}{Figura 99 --}

\renewcommand{\cfttabpresnum}{Tabela~}
\renewcommand{\cfttabaftersnum}{ --}
\settowidth{\cfttabnumwidth}{Tabela 99 --}

% Configurações para a Lista de Quadros
\newlength{\cftquadnumwidth}
\newcommand{\cftquadpresnum}{Quadro~}
\newcommand{\cftquadaftersnum}{ --}
\settowidth{\cftquadnumwidth}{\cftquadpresnum 99\cftquadaftersnum}

% Configurações para a Lista de Gráficos
\newlength{\cftgrafnumwidth}
\newcommand{\cftgrafpresnum}{Gráfico~}
\newcommand{\cftgrafaftersnum}{ --}
\settowidth{\cftgrafnumwidth}{\cftgrafpresnum 99\cftgrafaftersnum}

% Atalhos de referência
\newcommand{\tabref}[1]{Tabela~\ref{#1}}
\newcommand{\figref}[1]{Figura~\ref{#1}}
\newcommand{\secref}[1]{Seção~\ref{#1}}
\newcommand{\chapref}[1]{Capítulo~\ref{#1}}
\newcommand{\eqnref}[1]{Equação~\eqref{#1}}

% Recuos de seção
\setlength{\cftsecindent}{0pt}
\setlength{\cftsubsecindent}{0pt}

% Formatação do título do sumário e listas
\renewcommand{\cfttoctitlefont}{\hfil \normalfont \centering \bfseries\MakeUppercase}
\renewcommand{\cftaftertoctitle}{\\[\baselineskip]\mbox{}\hfill{\normalfont Página}}
\setlength{\cftbeforetoctitleskip}{-1.5em}
\setlength{\cftaftertoctitleskip}{0pt}

%-------------------------------------------------------
% AMBIENTES DE TEOREMAS E APLICAÇÕES
%-------------------------------------------------------
\newtheorem{theorem}{Teorema}[chapter]
\newtheorem{teorema}{Teorema}[chapter]
\newtheorem{corollary}{Corolário}[theorem]
\newtheorem{corolario}{Corolário}[chapter]
\newtheorem{lemma}[theorem]{Lema}

\theoremstyle{remark}
\newtheorem{remark}{Observação}[chapter]

\theoremstyle{definition}
\newtheorem{definition}{Definição}[chapter]
\newtheorem{exemplo}{Exemplo}[chapter]

\theoremstyle{plain}
\newtheorem{lema}{Lema}[chapter]
\newtheorem{prop}{Proposição}[chapter]
\newtheorem{cor}{Corolário}[chapter]

\declaretheorem[style=definition,qed=$\filledmedtriangleleft$, name=Aplicação, numberwithin=chapter]{ap}

\theoremstyle{definition}
\newtheorem*{nota}{Notação}

%-------------------------------------------------------
% INJEÇÃO DOS PREFIXOS NAS LISTAS (Gráfico e Quadro)
%-------------------------------------------------------
\makeatletter
% Formatação da Lista de Gráficos
\renewcommand{\l@grafico}[2]{%
  \begingroup
  \renewcommand{\numberline}[1]{\hb@xt@\@tempdima{\cftgrafpresnum##1\cftgrafaftersnum\hfil}}%
  \@dottedtocline{1}{0em}{\cftgrafnumwidth}{#1}{#2}%
  \endgroup
}

% Formatação da Lista de Quadros
\renewcommand{\l@quadro}[2]{%
  \begingroup
  \renewcommand{\numberline}[1]{\hb@xt@\@tempdima{\cftquadpresnum##1\cftquadaftersnum\hfil}}%
  \@dottedtocline{1}{0em}{\cftquadnumwidth}{#1}{#2}%
  \endgroup
}

% Comandos nativos para geração de listas pelo sumário global
\renewcommand{\listoffigures}{\@starttoc{lof}}
\renewcommand{\listoftables}{\@starttoc{lot}}
\makeatother

%-------------------------------------------------------
% REFERÊNCIAS CRUZADAS E LINKS
%-------------------------------------------------------
\usepackage[
    colorlinks=true, % Ativa cores nos links
    linkcolor=black, % Cor dos links internos
    citecolor=black, % Cor dos links das citações bibliográficas
    urlcolor=black   % Cor dos links de internet
]{hyperref}

%-------------------------------------------------------
% CITAÇÕES ABNT E BACKREF
%-------------------------------------------------------
\usepackage[brazilian,hyperpageref]{backref}
\usepackage[alf]{abntex2cite}

% Configurações do pacote backref
\renewcommand{\backrefpagesname}{Citada na(s) página(s):~}
\renewcommand{\backref}{}
\renewcommand*{\backrefalt}[4]{
    \ifcase #1 %
        Nenhuma citação no texto.%
    \or
        Citação na pág. #2.%
    \else
        Citação nas págs. #2.%
    \fi}

%Adicionados por mim
\usepackage{array}
\usepackage{tabularx}